\documentclass[a4paper, 12pt]{article}
\usepackage[utf8]{inputenc}
\usepackage[T1]{fontenc}
\usepackage[portuguese]{babel}
\usepackage{amsmath}
\usepackage{amssymb}
\usepackage{hyperref}
\usepackage{abstract}
\usepackage{times}
\usepackage{setspace}
\usepackage[a4paper, top=3cm, bottom=2cm, left=3cm, right=2cm]{geometry}
\usepackage{titlesec}

\onehalfspacing 
\setlength{\parindent}{1.25cm}

\titleformat{\section}
  {\normalfont\normalsize\bfseries}
  {\thesection}
  {0.5em}
  {\uppercase}
\titlespacing*{\section}{0pt}{1.5cm}{0.5cm}


\setlength{\absleftindent}{0cm}
\setlength{\absrightindent}{0cm}

\begin{document}

\thispagestyle{empty} 

\begin{center}
\vspace*{3cm}

\Large\textbf{\uppercase{EconoQuiz: Uma Plataforma Gamificada para a Promoção da Literacia Financeira e o Suporte à ODS 8 (Trabalho Decente e Crescimento Econômico)}}
\par
\vspace{2cm}

\normalsize
\begin{spacing}{1.0}
\begin{tabular}{c}
Abnoan Gabriel$^{1}$\\
José Danilo$^{2}$\\
Jéssica Isabela$^{3}$
\end{tabular}
\end{spacing}
\end{center}

\begin{spacing}{1.0}
\footnotetext[1]{
    \scriptsize UFERSA - UNIVERSIDADE FEDERAL DO SEMI-ÁRIDO.
}
\end{spacing}

\newpage
\thispagestyle{plain} 


\begin{center}
\normalfont\normalsize\textbf{RESUMO}
\end{center}

\begin{spacing}{1.0}
\noindent Este artigo apresenta o desenvolvimento e a arquitetura do \textbf{EconoQuiz}, uma aplicação de quiz gamificada projetada para aumentar a literacia financeira de forma acessível e envolvente. O projeto se alinha com a Meta 8.10 do \textbf{Objetivo de Desenvolvimento Sustentável (ODS) 8} (Trabalho Decente e Crescimento Econômico), buscando fortalecer a inclusão financeira. Detalhamos os requisitos funcionais e não funcionais que garantem uma experiência de aprendizado robusta, com partidas baseadas em níveis de dificuldade e pontuações que atualizam um ranking em tempo real. O EconoQuiz é proposto como uma solução tecnológica escalável para o desenvolvimento de habilidades econômicas essenciais.
\end{spacing}

\vspace{1cm}
\begin{spacing}{1.0}
\noindent \textbf{Palavras-chave:} Literacia Financeira; Gamificação; ODS 8; Desenvolvimento Sustentável; Tecnologia Educacional.
\end{spacing}


\newpage
\begin{spacing}{1.5} 

\section{Introdução}

A literacia financeira é um pilar essencial para a autonomia e o empoderamento econômico na sociedade contemporânea. A ausência de habilidades básicas em orçamento, gestão de dívidas e planejamento de investimentos restringe o acesso a serviços financeiros e limita as oportunidades de crescimento individual, o que impacta diretamente a capacidade produtiva e o desenvolvimento social. Em resposta a esse desafio, o presente trabalho descreve o desenvolvimento e a arquitetura do \textbf{EconoQuiz}, uma aplicação móvel e web desenhada para transformar a educação financeira em uma experiência envolvente e acessível.

O projeto EconoQuiz está intrinsecamente ligado ao \textbf{Objetivo de Desenvolvimento Sustentável (ODS) 8: Trabalho Decente e Crescimento Econômico} da Organização das Nações Unidas (ONU), especificamente à \textbf{Meta 8.10}, que visa fortalecer a inclusão financeira. Ao fornecer uma ferramenta de aprendizado prático e motivador, o EconoQuiz contribui para a capacitação de jovens e adultos com as habilidades financeiras necessárias para a participação plena no mercado de trabalho e na economia.

O EconoQuiz é estruturado como um sistema de quiz gamificado. As partidas são baseadas em questões com níveis de dificuldade variados, atribuindo pontuações que determinam o \textbf{nível de jogo} e a posição do usuário em um \textbf{ranking atualizado em tempo real}. O sistema é suportado por um conjunto robusto de requisitos funcionais e não funcionais rigorosos, garantindo uma plataforma escalável e de alto desempenho. Este artigo detalhará a relevância, o design e a implementação dessa solução como um agente de mudança social e econômica.

\end{spacing}
\end{document}